%%%%%%%%%%%%%%%%
%%% lev 26 %%%
%%%%%%%%%%%%%%%%

\Chapter{26}

\Verse{1}
{
	\hP{לֹא תַעֲשׂוּ לָכֶם אֱלִילִם וּפֶסֶל וּמַצֵּבָה לֹא תָקִימוּ לָכֶם וְאֶבֶן מַשְׂכִּית לֹא תִתְּנוּ בְּאַרְצְכֶם לְהִשְׁתַּחֲוֺת עָלֶיהָ כִּי אֲנִי יְהֹוָה אֱלֹהֵיכֶם׃}
}
{
	\eP{
		“You shall not make idols, and you shall not set up an
		image or a pillar, and you shall not put a carved stone in
		your land to bow at it, because I am YHWH, your {\God}.
	}
}
{
}

\Verse{2}
{
	\hP{אֶת שַׁבְּתֹתַי תִּשְׁמֹרוּ וּמִקְדָּשִׁי תִּירָאוּ אֲנִי יְהֹוָה׃}
}
{
	\eP{
		“You shall observe my Sabbaths and fear my sanctuary. I am
		YHWH.
	}
}
{
%	\footnote{
%		observe my Sabbaths and fear my sanctuary. What is the reason for
%		mentioning these two here together: Sabbaths and sanctuary? This
%		pronouncement brings together the sanctification of time and the
%		sanctification of space. The Sabbath is the most sacred time. The
%		sanctuary (meaning first the Tabernacle and later the Temple) is the
%		most sacred space. This dual command thus embraces everything.
%	}
}

\Verse{3}
{
	\hP{אִם בְּחֻקֹּתַי תֵּלֵכוּ וְאֶת מִצְוֺתַי תִּשְׁמְרוּ וַעֲשִׂיתֶם אֹתָם׃}
}
{
	\eP{
		BY MY LAWS “If you will go by my laws, and if you will
		observe my commandments, and you will do them:
	}
}
{
%	\footnote{
%		If you will go by my laws. The key to the relationship between
%		Leviticus’s laws and the Bible’s story is the list of blessings and
%		curses here in Leviticus 26. The list comes at the conclusion and is
%		arguably the culmination of the book. It indicates that all of this law
%		has implications. Leviticus has definite connections to what has come
%		before it. The blessings-and-curses list is the connection to what will
%		follow. It says, plainly enough: the people’s fate will depend on
%		whether they follow these laws or not. The blessings include rain,
%		produce, satiety, security, peace, victory, population, divine presence,
%		and divine relationship. The curses include illness, defeat, wild
%		animals, pestilence, famine, divine rejection, destruction, dispersion,
%		and exile. The blessings and curses thus provide for the eventuality of
%		a variety of essential matters in the realms of nature, human relations
%		(especially political relations), and divine-human relationship. All
%		future stories of the fortunes of the people of Israel in these realms
%		are to be understood, according to this text, as being the results of
%		whether “you will go by my laws, and if you will observe my
%		commandments, and you will do them” (26:3) or “if you will reject my
%		laws, and if your souls will scorn my judgments so as not to do all my
%		commandments, so that you break my covenant” (26:15). Underscoring the
%		literary force of these promises and threats is not only their placement
%		at the conclusion of the book but the fact that the list of curses is
%		the most artfully written section in the book. I call it the most
%		artfully written on the basis of its wording as much as its content,
%		including words of irony (“you will flee though no one is pursuing you,”
%		26:17), words of horror (“you will eat your sons’ flesh, and your
%		daughters’ flesh you will eat,” 26:29), and words of pathos (“the sound
%		of a driven leaf will chase them,” 26:36). The commandments thus are not
%		presented as a loosely relevant list. They are woven into the fabric of
%		the narrative as essential to the life of the community.
%	}
}

\Verse{4}
{
	\hP{וְנָתַתִּי גִשְׁמֵיכֶם בְּעִתָּם וְנָתְנָה הָאָרֶץ יְבוּלָהּ וְעֵץ הַשָּׂדֶה יִתֵּן פִּרְיוֹ׃}
}
{
	\eP{
		then I shall give your rains in their time, and the earth
		will give its crop, and the tree of the field will give
		its fruit,
	}
}
{
}

\Verse{5}
{
	\hP{וְהִשִּׂיג לָכֶם דַּיִשׁ אֶת בָּצִיר וּבָצִיר יַשִּׂיג אֶת זָרַע וַאֲכַלְתֶּם לַחְמְכֶם לָשֹׂבַע וִישַׁבְתֶּם לָבֶטַח בְּאַרְצְכֶם׃}
}
{
	\eP{
		and threshing will extend to vintage for you, and vintage
		will extend to seeding, and you will eat your bread to the
		full, and you will live in security in your land,
	}
}
{
}

\Verse{6}
{
	\hP{וְנָתַתִּי שָׁלוֹם בָּאָרֶץ וּשְׁכַבְתֶּם וְאֵין מַחֲרִיד וְהִשְׁבַּתִּי חַיָּה רָעָה מִן הָאָרֶץ וְחֶרֶב לֹא תַעֲבֹר בְּאַרְצְכֶם׃}
}
{
	\eP{
		and I shall give peace in the land, and you will lie down
		with no one making you afraid, and I shall make wild
		animals cease from the land, and a sword will not pass
		through your land,
	}
}
{
}

\Verse{7}
{
	\hP{וּרְדַפְתֶּם אֶת אֹיְבֵיכֶם וְנָפְלוּ לִפְנֵיכֶם לֶחָרֶב׃}
}
{
	\eP{
		and you will chase your enemies, and they will fall in
		front of you by the sword,
	}
}
{
}

\Verse{8}
{
	\hP{וְרָדְפוּ מִכֶּם חֲמִשָּׁה מֵאָה וּמֵאָה מִכֶּם רְבָבָה יִרְדֹּפוּ וְנָפְלוּ אֹיְבֵיכֶם לִפְנֵיכֶם לֶחָרֶב׃}
}
{
	\eP{
		and five of you will chase a hundred, and a hundred of you
		will chase ten thousand, and your enemies will fall in
		front of you by the sword,
	}
}
{
}

\Verse{9}
{
	\hP{וּפָנִיתִי אֲלֵיכֶם וְהִפְרֵיתִי אֶתְכֶם וְהִרְבֵּיתִי אֶתְכֶם וַהֲקִימֹתִי אֶת בְּרִיתִי אִתְּכֶם׃}
}
{
	\eP{
		and I shall turn to you, and I shall make you fruitful and
		make you multiply, and I shall establish my covenant with
		you,
	}
}
{
}

\Verse{10}
{
	\hP{וַאֲכַלְתֶּם יָשָׁן נוֹשָׁן וְיָשָׁן מִפְּנֵי חָדָשׁ תּוֹצִיאוּ׃}
}
{
	\eP{
		and you will eat the oldest of the old, and you will take
		out the old because of the new,
	}
}
{
%	\footnote{
%		the oldest of the old. Old, stored produce will still make good food.
%		Footnote 26:10. take out the old because of the new. One will never run
%		out of food. There will be good stored produce right up to the time that
%		the new crops come in, and the old will have to be taken out in order to
%		make room for the new.
%	}
}

\Verse{11}
{
	\hP{וְנָתַתִּי מִשְׁכָּנִי בְּתוֹכְכֶם וְלֹא תִגְעַל נַפְשִׁי אֶתְכֶם׃}
}
{
	\eP{
		and I shall put my Tabernacle among you, and my soul will
		not scorn you,
	}
}
{
%	\footnote{
%		my soul. Does God have a soul? This is the only place in the Torah that
%		refers to the deity’s nepeš (although a related verbal form, wayyinnpaš,
%		occurs in Exod 31:17). Elsewhere, the word refers to the living quality
%		in humans and animals and is associated with breath. It is usually
%		understood to mean soul, person, being, and life. It might possibly help
%		us to understand what is meant by creation in the image of God, but that
%		seems unlikely since animals are said to have a nepeš there as well, but
%		they are not said to be in the divine image (Gen 1:24–27). We must be
%		cautious in using the word’s occurrences here to conclude anything about
%		the Torah’s conception of God, because both of these occurrences are in
%		the phrase “my soul will scorn.” This phrase may simply have been a
%		known expression. It also occurs here with humans as the subject
%		(26:15); and it occurs in Jeremiah as well (14:19). It may refer to the
%		personal aspect of God, or it may rather refer to God in no specific
%		way. Rashi (following Siphra) takes the phrase to refer to the departure
%		of the divine presence (Shechinah). Here, too, we must be cautious,
%		first, because that is also the understanding of other metaphorical
%		phrases in the Tanak. Thus the Targum takes the phrase “God hides His
%		face” to mean the departure of the Shechinah as well. And, second, the
%		word Shechinah never occurs in the Torah or anywhere in the Hebrew
%		Bible.
%	}
}

\Verse{12}
{
	\hP{וְהִתְהַלַּכְתִּי בְּתוֹכְכֶם וְהָיִיתִי לָכֶם לֵאלֹהִים וְאַתֶּם תִּהְיוּ לִי לְעָם׃}
}
{
	\eP{
		and I shall walk among you, and I shall be {\God} to you, and
		you will be a people to me.
	}
}
{
%	\footnote{
%		walk among you. This same term is used to describe God’s walking among
%		humans in the garden of Eden (the Hebrew root in the Hitpael; Gen 3:8).
%		This blessing thus hints at a return of the closeness to the divine
%		presence that existed prior to the initial divine-human conflict in
%		Eden. Footnote 26:12. I shall be God to you, and you will be a people to
%		me. The last blessing of the list echoes what is thought to be the
%		wording of the marriage ceremony in ancient Israel: “I shall be a
%		husband to you, and you will be a wife to me”—as well as the adoption
%		ceremony: “I shall be a parent to you, and you will be a son/daughter to
%		me” (see, e.g., 2 Sam 7:14).
%	}
}

\Verse{13}
{
	\hP{אֲנִי יְהֹוָה אֱלֹהֵיכֶם אֲשֶׁר הוֹצֵאתִי אֶתְכֶם מֵאֶרֶץ מִצְרַיִם מִהְיֹת לָהֶם עֲבָדִים וָאֶשְׁבֹּר מֹטֹת עֻלְּכֶם וָאוֹלֵךְ אֶתְכֶם קוֹמְמִיּוּת׃}
}
{
	\eP{
		I am YHWH, your {\God}, who brought you out from the land of
		Egypt, from being slaves to them, and I broke the beams of
		your yoke and had you go standing tall.
	}
}
{
}

\Verse{14}
{
	\hP{וְאִם לֹא תִשְׁמְעוּ לִי וְלֹא תַעֲשׂוּ אֵת כּל הַמִּצְוֺת הָאֵלֶּה׃}
}
{
	\eP{
		“But if you will not listen to me and not do all these
		commandments,
	}
}
{
}

\Verse{15}
{
	\hP{וְאִם בְּחֻקֹּתַי תִּמְאָסוּ וְאִם אֶת מִשְׁפָּטַי תִּגְעַל נַפְשְׁכֶם לְבִלְתִּי עֲשׂוֹת אֶת כּל מִצְוֺתַי לְהַפְרְכֶם אֶת בְּרִיתִי׃}
}
{
	\eP{
		and if you will reject my laws, and if your souls will
		scorn my judgments so as not to do all my commandments, so
		that you break my covenant,
	}
}
{
}

\Verse{16}
{
	\hP{אַף אֲנִי אֶעֱשֶׂה זֹּאת לָכֶם וְהִפְקַדְתִּי עֲלֵיכֶם בֶּהָלָה אֶת הַשַּׁחֶפֶת וְאֶת הַקַּדַּחַת מְכַלּוֹת עֵינַיִם וּמְדִיבֹת נָפֶשׁ וּזְרַעְתֶּם לָרִיק זַרְעֲכֶם וַאֲכָלֻהוּ אֹיְבֵיכֶם׃}
}
{
	\eP{
		I, too, I shall do this to you: Then I shall appoint
		terror over you with consumption and with fever,
		exhausting eyes and grieving the soul, and you will sow
		your seed in vain, for your enemies will eat it,
	}
}
{
}

\Verse{17}
{
	\hP{וְנָתַתִּי פָנַי בָּכֶם וְנִגַּפְתֶּם לִפְנֵי אֹיְבֵיכֶם וְרָדוּ בָכֶם שֹׂנְאֵיכֶם וְנַסְתֶּם וְאֵין רֹדֵף אֶתְכֶם׃}
}
{
	\eP{
		and I shall set my face against you, and you will be
		struck before your enemies, and ones who hate you will
		dominate you, and you will flee though no one is pursuing
		you.
	}
}
{
}

\Verse{18}
{
	\hP{וְאִם עַד אֵלֶּה לֹא תִשְׁמְעוּ לִי וְיָסַפְתִּי לְיַסְּרָה אֶתְכֶם שֶׁבַע עַל חַטֹּאתֵיכֶם׃}
}
{
	\eP{
		And if, even after these, you will not listen to me, then
		I shall add seven times more to discipline you for your
		sins:
	}
}
{
}

\Verse{19}
{
	\hP{וְשָׁבַרְתִּי אֶת גְּאוֹן עֻזְּכֶם וְנָתַתִּי אֶת שְׁמֵיכֶם כַּבַּרְזֶל וְאֶת אַרְצְכֶם כַּנְּחֻשָׁה׃}
}
{
	\eP{
		and I shall break your strong pride, and I shall make your
		skies like iron and your land like bronze,
	}
}
{
}

\Verse{20}
{
	\hP{וְתַם לָרִיק כֹּחֲכֶם וְלֹא תִתֵּן אַרְצְכֶם אֶת יְבוּלָהּ וְעֵץ הָאָרֶץ לֹא יִתֵּן פִּרְיוֹ׃}
}
{
	\eP{
		and your power will be used up in vain, and your land will
		not give its crop, and the tree of the land will not give
		its fruit.
	}
}
{
}

\Verse{21}
{
	\hP{וְאִם תֵּלְכוּ עִמִּי קֶרִי וְלֹא תֹאבוּ לִשְׁמֹעַ לִי וְיָסַפְתִּי עֲלֵיכֶם מַכָּה שֶׁבַע כְּחַטֹּאתֵיכֶם׃}
}
{
	\eP{
		And if you will go on with me in defiance and will not be
		willing to listen to me, then I shall add seven times more
		striking on you, corresponding to your sins:
	}
}
{
}

\Verse{22}
{
	\hP{וְהִשְׁלַחְתִּי בָכֶם אֶת חַיַּת הַשָּׂדֶה וְשִׁכְּלָה אֶתְכֶם וְהִכְרִיתָה אֶת בְּהֶמְתְּכֶם וְהִמְעִיטָה אֶתְכֶם וְנָשַׁמּוּ דַּרְכֵיכֶם׃}
}
{
	\eP{
		and I shall let loose the wild animal among you, and it
		will bereave you, and it will cut off your cattle, and it
		will diminish you, and your roads will be desolated.
	}
}
{
}

\Verse{23}
{
	\hP{וְאִם בְּאֵלֶּה לֹא תִוָּסְרוּ לִי וַהֲלַכְתֶּם עִמִּי קֶרִי׃}
}
{
	\eP{
		And if you will not be disciplined for me by these, and
		you will go on with me in defiance,
	}
}
{
}

\Verse{24}
{
	\hP{וְהָלַכְתִּי אַף אֲנִי עִמָּכֶם בְּקֶרִי וְהִכֵּיתִי אֶתְכֶם גַּם אָנִי שֶׁבַע עַל חַטֹּאתֵיכֶם׃}
}
{
	\eP{
		then I shall go on with you, I also, in defiance. And I
		shall strike you, I too, seven times more for your sins:
	}
}
{
}

\Verse{25}
{
	\hP{וְהֵבֵאתִי עֲלֵיכֶם חֶרֶב נֹקֶמֶת נְקַם בְּרִית וְנֶאֱסַפְתֶּם אֶל עָרֵיכֶם וְשִׁלַּחְתִּי דֶבֶר בְּתוֹכְכֶם וְנִתַּתֶּם בְּיַד אוֹיֵב׃}
}
{
	\eP{
		and I shall bring a sword over you, requiting covenant
		retribution, and you will be gathered to your cities, and
		I shall let an epidemic go among you, and you will be put
		in your enemy’s hand.
	}
}
{
}

\Verse{26}
{
	\hP{בְּשִׁבְרִי לָכֶם מַטֵּה לֶחֶם וְאָפוּ עֶשֶׂר נָשִׁים לַחְמְכֶם בְּתַנּוּר אֶחָד וְהֵשִׁיבוּ לַחְמְכֶם בַּמִּשְׁקָל וַאֲכַלְתֶּם וְלֹא תִשְׂבָּעוּ׃}
}
{
	\eP{
		When I break the staff of bread for you, then ten women
		will bake your bread in one oven, and they will give back
		your bread by weight, and you will eat and not be full.
	}
}
{
}

\Verse{27}
{
	\hP{וְאִם בְּזֹאת לֹא תִשְׁמְעוּ לִי וַהֲלַכְתֶּם עִמִּי בְּקֶרִי׃}
}
{
	\eP{
		And if, through this, you will not listen to me, and you
		will go on with me in defiance,
	}
}
{
}

\Verse{28}
{
	\hP{וְהָלַכְתִּי עִמָּכֶם בַּחֲמַת קֶרִי וְיִסַּרְתִּי אֶתְכֶם אַף אָנִי שֶׁבַע עַל חַטֹּאתֵיכֶם׃}
}
{
	\eP{
		then I shall go on with you in a fury of defiance, and I
		shall discipline you, I also, seven times more for your
		sins:
	}
}
{
}

\Verse{29}
{
	\hP{וַאֲכַלְתֶּם בְּשַׂר בְּנֵיכֶם וּבְשַׂר בְּנֹתֵיכֶם תֹּאכֵלוּ׃}
}
{
	\eP{
		and you will eat your sons’ flesh, and your daughters’
		flesh you will eat,
	}
}
{
%	\footnote{
%		you will eat your sons’ flesh and your daughters’ flesh. The blessings
%		and curses are not rewards and punishments. They are a formal part of
%		the biblical covenants (Leviticus 26; Deuteronomy 28)—with exact
%		parallels in legal contracts of the ancient Near East. They express the
%		outcomes of fulfillment or nonfulfillment of the covenant’s terms. A
%		people who is faithful to its God and keeps the Sabbath and honors its
%		parents and does not steal will be prosperous and enduring in its land.
%		A people who loses sight of its commitments and values will suffer. (It
%		should not be hard to think of contemporary parallels.) The frightful
%		curse of eating one’s own children comes true centuries later in one of
%		the most horrifying stories in the Bible (2 Kings 6:24–30). It is not
%		presented there as a punishment. It rather conveys the terrible state of
%		things in Israel in the wake of the heretical reigns of Kings Ahab and
%		Jehoram. There is a big difference between a punishment and a curse,
%		between a threat and a warning. This explanation is still distressing,
%		but it does not so easily picture the deity as malicious, as people have
%		sometimes taken these curses to mean. The God of the Hebrew Bible is not
%		the “Old Testament God of Wrath,” but rather a deity who is torn between
%		mercy and justice, between affection for humankind and regret over the
%		continuous conflict with them. The curses are a sad outcome of a certain
%		kind of human behavior. But, then, the blessings are the outcome of the
%		other kind.
%	}
}

\Verse{30}
{
	\hP{וְהִשְׁמַדְתִּי אֶת בָּמֹתֵיכֶם וְהִכְרַתִּי אֶת חַמָּנֵיכֶם וְנָתַתִּי אֶת פִּגְרֵיכֶם עַל פִּגְרֵי גִּלּוּלֵיכֶם וְגָעֲלָה נַפְשִׁי אֶתְכֶם׃}
}
{
	\eP{
		and I shall destroy your high places, and I shall cut off
		your incense altars, and I shall put your carcasses on
		your idols’ carcasses, and my soul will scorn you,
	}
}
{
%	\footnote{
%		high places. These were large stone platforms where sacrifices were
%		performed. Some were used for worship of YHWH, and others were used for
%		worship of pagan gods; but both kinds are regularly forbidden in the
%		Tanak. High places to the God of Israel are regarded as violations of
%		the laws of centralization of worship at a single site (see comments on
%		Lev 1:3 and 17:4). Footnote 26:30. incense altars. These were stands,
%		usually designed, on which incense was burned. Both high places and
%		incense altars have been excavated in Israel.
%	}
}

\Verse{31}
{
	\hP{וְנָתַתִּי אֶת עָרֵיכֶם חרְבָּה וַהֲשִׁמּוֹתִי אֶת מִקְדְּשֵׁיכֶם וְלֹא אָרִיחַ בְּרֵיחַ נִיחֹחֲכֶם׃}
}
{
	\eP{
		and I shall make your cities a ruin, and I shall devastate
		your holy places, and I shall not smell your pleasant
		smell,
	}
}
{
%	\footnote{
%		not smell your pleasant smell. The “pleasant smell” has been referred to
%		many times as the product of the burning incense that accompanies the
%		sacrifices. Starting with Noah’s sacrifice after the flood, the smell of
%		sacrifices is meant to be attractive to the deity, but the curse here is
%		that God will refuse to smell it.
%	}
}

\Verse{32}
{
	\hP{וַהֲשִׁמֹּתִי אֲנִי אֶת הָאָרֶץ וְשָׁמְמוּ עָלֶיהָ אֹיְבֵיכֶם הַיֹּשְׁבִים בָּהּ׃}
}
{
	\eP{
		and I, I, shall devastate the land, so your enemies who
		live in it will be astonished,
	}
}
{
}

\Verse{33}
{
	\hP{וְאֶתְכֶם אֱזָרֶה בַגּוֹיִם וַהֲרִיקֹתִי אַחֲרֵיכֶם חָרֶב וְהָיְתָה אַרְצְכֶם שְׁמָמָה וְעָרֵיכֶם יִהְיוּ חרְבָּה׃}
}
{
	\eP{
		and you: I shall scatter among the nations, and I shall
		unsheathe my sword after you, and your land will be a
		devastation, and your cities will be a ruin.
	}
}
{
%	\footnote{
%		scatter among the nations. Among the worst curses was the curse of
%		exile: for a people to be driven off their land and scattered around the
%		world. This possibility was a reality in the ancient Near East, where
%		exile was a horrifying fate. This curse would come to pass for Israel
%		later in the Tanak, when they would come to know exile at the hands of
%		the Assyrians; and it would come to pass for Judah still later at the
%		hands of the Babylonians. But they would also come to know the blessing
%		of return. Footnote 26:33. unsheathe my sword. To appreciate the horror
%		of the metaphor of God saying He will go after Israel unsheathing a
%		sword, one should recall that this is exactly what the Egyptians say
%		that they will do to Israel in the Song of the Sea (Exod 15:9). That is,
%		God, who is Israel’s protector, threatens a curse in which God would
%		become like Israel’s enemies. It is one’s worst nightmare: that one’s
%		loving, protecting parents become a horrible threat.
%	}
}

\Verse{34}
{
	\hP{אָז תִּרְצֶה הָאָרֶץ אֶת שַׁבְּתֹתֶיהָ כֹּל יְמֵי השַּׁמָּה וְאַתֶּם בְּאֶרֶץ אֹיְבֵיכֶם אָז תִּשְׁבַּת הָאָרֶץ וְהִרְצָת אֶת שַׁבְּתֹתֶיהָ׃}
}
{
	\eP{
		Then the land will accept its Sabbaths: all the days of
		devastation, while you are in your enemies’ land. Then the
		land will cease and accept its Sabbaths.
	}
}
{
}

\Verse{35}
{
	\hP{כּל יְמֵי השַּׁמָּה תִּשְׁבֹּת אֵת אֲשֶׁר לֹא שָׁבְתָה בְּשַׁבְּתֹתֵיכֶם בְּשִׁבְתְּכֶם עָלֶיהָ׃}
}
{
	\eP{
		All the days of the desolation it will cease, the amount
		that it did not cease on your Sabbaths when you lived on
		it.
	}
}
{
}

\Verse{36}
{
	\hP{וְהַנִּשְׁאָרִים בָּכֶם וְהֵבֵאתִי מֹרֶךְ בִּלְבָבָם בְּאַרְצֹת אֹיְבֵיהֶם וְרָדַף אֹתָם קוֹל עָלֶה נִדָּף וְנָסוּ מְנֻסַת חֶרֶב וְנָפְלוּ וְאֵין רֹדֵף׃}
}
{
	\eP{
		And those of you who remain: I shall bring faintness in
		their hearts in their enemies’ lands, and the sound of a
		driven leaf will chase them, and they will flee like the
		flight from a sword, and they will fall when there is no
		one pursuing,
	}
}
{
}

\Verse{37}
{
	\hP{וְכָשְׁלוּ אִישׁ בְּאָחִיו כְּמִפְּנֵי חֶרֶב וְרֹדֵף אָיִן וְלֹא תִהְיֶה לָכֶם תְּקוּמָה לִפְנֵי אֹיְבֵיכֶם׃}
}
{
	\eP{
		and they will stumble, each over his brother, as in front
		of a sword, when there is no one pursuing, and you will
		not have a footing in front of your enemies,
	}
}
{
}

\Verse{38}
{
	\hP{וַאֲבַדְתֶּם בַּגּוֹיִם וְאָכְלָה אֶתְכֶם אֶרֶץ אֹיְבֵיכֶם׃}
}
{
	\eP{
		and you will perish among the nations, and your enemies’
		land will eat you up,
	}
}
{
}

\Verse{39}
{
	\hR{וְהַנִּשְׁאָרִים בָּכֶם יִמַּקּוּ בַּעֲוֺנָם בְּאַרְצֹת אֹיְבֵיכֶם וְאַף בַּעֲוֺנֹת אֲבֹתָם אִתָּם יִמָּקּוּ׃}
}
{
	\eR{
		and those of you who remain will rot for their crime in
		your enemies’ lands, and also for their fathers’ crimes
		with them they will rot.
	}
}
{
}

\Verse{40}
{
	\hR{וְהִתְוַדּוּ אֶת עֲוֺנָם וְאֶת עֲוֺן אֲבֹתָם בְּמַעֲלָם אֲשֶׁר מָעֲלוּ בִי וְאַף אֲשֶׁר הָלְכוּ עִמִּי בְּקֶרִי׃}
}
{
	\eR{
		“When they will confess their crime and their fathers’
		crime for their breach that they made against me, and also
		that they went on with me in defiance,
	}
}
{
%	\footnote{
%		their crime and their fathers’ crime. On one hand, this is yet another
%		reminder that it takes generations before the God of Israel will take
%		such actions. On the other hand, it means that the Israelites must
%		recognize that it is not only their own offenses that have led to these
%		things. They must come to terms with their parents’ offenses as well.
%		The ultimate revelation of the divine at Sinai includes the fact that
%		the deity “reckons fathers’ crime on children and on children’s
%		children, on third generations and on fourth generations.” Now the
%		perspective is turned to those children, generations later; and we learn
%		from the Torah, no less than from Freud, that humans must look back not
%		only at the good that their parents and grandparents did but also at
%		their errors and faults. It is an essential part of understanding how we
%		came to be what we are.
%	}
}

\Verse{41}
{
	\hR{אַף אֲנִי אֵלֵךְ עִמָּם בְּקֶרִי וְהֵבֵאתִי אֹתָם בְּאֶרֶץ אֹיְבֵיהֶם אוֹ אָז יִכָּנַע לְבָבָם הֶעָרֵל וְאָז יִרְצוּ אֶת עֲוֺנָם׃}
}
{
	\eR{
		so I, I also, would go on with them in defiance, and I
		brought them into their enemies’ land; or if then their
		uncircumcised heart will be humbled and then they will
		accept their crime,
	}
}
{
}

\Verse{42}
{
	\hR{וְזָכַרְתִּי אֶת בְּרִיתִי יַעֲקוֹב וְאַף אֶת בְּרִיתִי יִצְחָק וְאַף אֶת בְּרִיתִי אַבְרָהָם אֶזְכֹּר וְהָאָרֶץ אֶזְכֹּר׃}
}
{
	\eR{
		then I shall remember my covenant of Jacob, and I shall
		remember also my covenant of Isaac and also my covenant of
		Abraham, and I shall remember the land.
	}
}
{
}

\Verse{43}
{
	\hR{וְהָאָרֶץ תֵּעָזֵב מֵהֶם וְתִרֶץ אֶת שַׁבְּתֹתֶיהָ בּהְשַׁמָּה מֵהֶם וְהֵם יִרְצוּ אֶת עֲוֺנָם יַעַן וּבְיַעַן בְּמִשְׁפָּטַי מָאָסוּ וְאֶת חֻקֹּתַי גָּעֲלָה נַפְשָׁם׃}
}
{
	\eR{
		And the land will be left from them, and it will accept
		its Sabbaths while it is devastated from them, and they
		will accept their crime, because and only because they
		rejected my judgments and their soul scorned my laws.
	}
}
{
}

\Verse{44}
{
	\hR{וְאַף גַּם זֹאת בִּהְיוֹתָם בְּאֶרֶץ אֹיְבֵיהֶם לֹא מְאַסְתִּים וְלֹא גְעַלְתִּים לְכַלֹּתָם לְהָפֵר בְּרִיתִי אִתָּם כִּי אֲנִי יְהֹוָה אֱלֹהֵיהֶם׃}
}
{
	\eR{
		“And even despite this, when they were in their enemies’
		land I did not reject them and did not scorn them, to
		finish them, to break my covenant with them, because I am
		YHWH, their {\God}!
	}
}
{
}

\Verse{45}
{
	\hR{וְזָכַרְתִּי לָהֶם בְּרִית רִאשֹׁנִים אֲשֶׁר הוֹצֵאתִי אֹתָם מֵאֶרֶץ מִצְרַיִם לְעֵינֵי הַגּוֹיִם לִהְיוֹת לָהֶם לֵאלֹהִים אֲנִי יְהֹוָה׃}
}
{
	\eR{
		So I shall remember for them the covenant of the first
		ones, whom I brought out from the land of Egypt before the
		eyes of the nations to be {\God} to them. I am YHWH.”
	}
}
{
}

\Verse{46}
{
	\hP{אֵלֶּה הַחֻקִּים וְהַמִּשְׁפָּטִים וְהַתּוֹרֹת אֲשֶׁר נָתַן יְהֹוָה בֵּינוֹ וּבֵין בְּנֵי יִשְׂרָאֵל בְּהַר סִינַי בְּיַד מֹשֶׁה׃}
}
{
	\eP{
		These are the laws and the judgments and the instructions
		that YHWH gave between Him and the children of Israel in
		Mount Sinai by Moses’ hand.
	}
}
{
%	\footnote{
%		These are the laws and the judgments and the instructions. The present
%		age plainly appears to value the ethical more than the ritual in
%		religion. Some even feel the need to justify ritual by attempting to
%		connect each ritual act to some ethical value: “We keep kosher to remind
%		us to care about animals; we wear fringes to remind us to be kind....”
%		This is misleading. Certainly ritual acts can have consequences in the
%		ethical realm, but that is not their reason for being. If we are to
%		understand Leviticus, we must have an appreciation of what ritual meant
%		in its society intrinsically. Ritual can have aesthetic, psychological,
%		and inspirational value. It has the potential of generating feelings of
%		closeness to the sacred as well as feelings of identity, stability, and
%		security. For some, it means, in a world of uncertainty, to be told what
%		the “right” things are, what acts one can perform to feel that one is
%		acting in accordance with nothing less than the will of the creator of
%		the universe. In the matter of purity, for example, one can believe
%		that, when one has had contact with any of a group of things that
%		intimidate us but that are physiological facts of life (blood and other
%		bodily fluids, bodily sores, death), one can do things that remove the
%		condition that this contact has produced—and thus remove the fear. In
%		the matter of sacrifice, as we have observed, the ritual has an impact
%		on one’s feeling of guilt. In the observance of the holidays, the ritual
%		contributes to feelings of festivity and communal celebration—and the
%		sanctification of time. And in the variety of elaborate ceremonies
%		performed at the Tent of Meeting, officiated by the specially garbed
%		priests, accompanied by the burning of incense, the ritual contributes
%		an aesthetic component to religion as well. All of this is
%		known—consciously or intuitively—to those who have found meaning and
%		even enjoyment in the practice of the rituals of their own religion. The
%		ritual and the ethical are two components of religion—and of
%		Leviticus—that do not justify each other, but rather unite and produce
%		mutual support. Indeed, it is instructive that Leviticus, a book that is
%		so fundamentally concerned with distinction, does not make any explicit
%		distinction between its ethical and its ritual laws. Sometimes they are
%		mixed together, but they are never identified as two distinct categories
%		of law. We do a disservice when we interpret ritual matters in terms of
%		ethical justifications. We move ourselves even farther from the
%		appreciation of ritual and the awe before the holy. Similarly, when we
%		say that the Hebrew word qdōš means a neutral kind of “separate” rather
%		than meaning “holy, sacred,” we further lose our feeling for holiness.
%		And perhaps we err similarly when we say that Hebrew t means “to miss
%		the mark” rather than “to sin,” to do something wrong that places us in
%		a state that needs to be remedied through atonement and/or sacrifice
%		and/or compensation.
%	}
}
